% Candidate concise methods; not included in the active manuscript.
\subsection{Controlled Adam Measurements and Timing}

Each of nine trajectories contains 1,024 full-objective Adam updates. The MLP
rates are $10^{-4}$, $3\cdot 10^{-4}$ and $10^{-3}$. ResNet-32 and the Transformer
use $3\cdot 10^{-4}$, $10^{-3}$ and $3\cdot 10^{-3}$.
All use $(\beta_1,\beta_2)=(0.9,0.999)$, $\epsilon=10^{-8}$, fresh moments,
FP32 training and FP64 diagnostic reductions, with TF32 disabled. There is
one frozen initialization per case, shared across its rates. There is no
clipping, decay, learning-rate schedule or Hessian measurement.

The MLP objective covers the frozen 5,000-image CIFAR-10 subset; ResNet-32
covers all 50,000 training images with 50 fixed ghost-batch-normalization
groups of 1,000; the Transformer covers all 2,088,600 prepared WikiText-2 target
tokens per update. Model, cache and initial tensor identities remain recorded
in Run 215. Parameter scope is all trainable parameters, with tied parameters
counted once. Zero-norm cosine and projected change are undefined.

In the figures, parameter state$t$ means$t$ completed updates. Loss and raw
gradient are measured at$\theta_t$; the displayed first moment has consumed
that gradient and is available before displacement$t+1$. In the manuscript's
update-index convention, these are$g_{s+1}$ and$m_{s+1}$ at state$s$. Integration
must use one notation consistently. No next-step observation is shifted into
the current state.

\subsection{Loss-Only Episodes and Earlier Crossings}

Search from state 64. At each unassigned state, use the minimum loss over the
preceding 16 states, excluding the current state. An episode starts at a rise
of at least 10\% above this reference. Freeze that reference and merge all
rebounds until the first subsequent loss strictly below it. Resume searching
after recovery. Retain unfinished recoveries at the final state. A large
episode has an observed peak at least 50\% above its baseline; two recoveries,
one in the highest-rate MLP and one in the highest-rate Transformer, are unfinished.

A qualifying earlier crossing goes from projected change at most two to above
two,1--8 states before episode onset. At that crossing, loss must remain less
than 5\% above the minimum over its own current and previous 16 states. Keep
every episode in the denominator, including misses. Lead time is onset minus
the earliest qualifying crossing in this limited window, summarized only
among episodes with a crossing. It is not measured relative to the loss peak.

The micro views select each highest-rate arm's completed episode with the
largest peak/baseline ratio, using loss alone. Each crop begins 10 states before
the peak and ends at the first subsequent loss below the crop-start value.
The crop reference and episode baseline can differ. Norms in this figure are
divided by their own crop-start value; they compare proportional changes,
not absolute gradient and momentum magnitudes.

The overview's non-loss prominent curves use complete 16-observation trailing
arithmetic means, right aligned without padding. Projected change is computed
pointwise before averaging; raw overlays are omitted. Its displayed
range is intentionally restricted to$[-0.1,2.2]$, without altering or clamping
the source values. The micro views, episode definitions and lead calculations
are unsmoothed. Repeated episodes are dependent, the design was motivated by
an inspected peak, and one initialization per case limits inference beyond
these trajectories. Counts and lead ranges are descriptive, not independent
replications or confidence intervals.
